\section{\(\lambda\)-definiowalność}
Funkcja \(f:\natural^k \to \natural\) jest \(\lambda\)-definiowalna przez term domknięty \(F\) jeśli dla dowolnych liczb \(n_1, \dots, n_k\)
\[ Fn_1n_2 \dots n_k =_\beta f(n_1,\dots,n_k) \]

\begin{theorem}[Church]
    Funkcje rekurencyjne są \(\lambda\)-definiowalne.
\end{theorem}
\begin{proof}
    Dowodzimy indukcją strukturalną po funkcjach rekurencyjnych:
    \begin{itemize}
        \item Funkcje inicjujące
        \begin{itemize}
        \singlespacing
            \item \(\lambda x_1\dots x_k.0\) -- zero
            \item \(\lambda nsz.s(nsz)\) -- następnik
            \item \(\lambda x_1 \dots x_k.x_i, \ i \leq k\) -- rzutowania
        \end{itemize}

        \item Domknięcie na złożenia
        \[ g \equiv \lambda x_1 \dots x_k.G(F_1x_1 \dots x_k) \dots(F_nx_1 \dots x_k) \]
    
        \item Domknięcie na rekursję prostą
        \begin{align*}
            M = \lambda F n x_1 \dots x_k.&\text{IF (iszero } n \text{) THEN } G x_1 \dots x_k \\
            &\text{ELSE } H \: (\text{pred } n) \: x_1\dots x_k \: (F (\text{pred } n) x_1 \dots x_k), \\
            f \equiv YM
        \end{align*}
        Dowód indukcyjny: \\
        \( YM \: 0 \: x_1 \dots x_k =_\beta M(YM) \: 0 \: x_1 \dots x_k =_\beta G x_1 \dots x_k =_\beta g(x_1, \dots, x_k) =_\beta f(0, x_1,\dots, x_k) \)
        \begin{align*}
            YM \: (n+1) \: x_1 \dots x_k &=_\beta M(YM) \: (n+1) \: x_1 \dots x_k \\
            &=_\beta H n x_1 \dots x_k ((YM) n x_1 \dots x_k) \\
            &=_\beta Hn x_1 \dots x_k (f(n, x_1, \dots, x_k)) \text{ (z zał. ind.)} \\
            &=_\beta h(n, x_1, \dots, x_k, f(n, x_1, \dots, x_k)) \\
            &=_\beta f(n, x_1, \dots, x_k)
        \end{align*}

        \item Domknięcie na \(\mu\)-rekursję
        \[
        \begin{aligned}
            M & = \lambda G y.\text{IF }(G y) \text{ THEN } y \text{ ELSE } G (y+1) \\[0.5em]
            f & \equiv YM
        \end{aligned}
    \]
    \end{itemize}
\end{proof}